% ============================================================
% CAS/Python ενότητες — Κεφάλαιο 5: Ιδιοτιμές & Ιδιοδιανύσματα
%
% QR code: qr_94ee53b0ae3e.png  (→ latex/qrcodes/)
% URL: https://nikmatz.github.io/ELADIC_GR/code/ch05/
% ============================================================

\casactivbox%
  {Maxima: Ιδιοτιμές — Χαρακτηριστικό Πολυώνυμο, Διαγωνοποίηση}%
  {%
    \label{cas:ch5_eigenvalues}
    Δίνεται ο πίνακας
    $A = \begin{pmatrix} 4 & 1 \\ 2 & 3 \end{pmatrix}$.
    \begin{enumerate}[label=(\alph*),itemsep=2pt]
      \item Υπολογίστε το χαρακτηριστικό πολυώνυμο
        $p(\lambda)=\det(A-\lambda I)$ με \texttt{charpoly}.
        Λύστε $p(\lambda)=0$ και επαληθεύστε $\operatorname{tr}(A)=\sum\lambda_i$,
        $\det(A)=\prod\lambda_i$.
      \item Για κάθε ιδιοτιμή βρείτε τον αντίστοιχο
        ιδιόχωρο (\texttt{nullspace}$(A-\lambda_i I)$).
        Επαληθεύστε $Av=\lambda v$.
      \item Για τον κατώτερο τριγωνικό πίνακα
        $C=\bigl[\begin{smallmatrix}2&0&0\\1&3&0\\0&1&2\end{smallmatrix}\bigr]$:
        υπολογίστε ιδιοτιμές με \texttt{eigenvectors(C)}.
        Γιατί οι ιδιοτιμές είναι τα διαγώνια στοιχεία;
      \item Κατασκευάστε $P=[v_1|v_2]$ και επαληθεύστε
        $A=PDP^{-1}$ (\texttt{invert}).
    \end{enumerate}
    \smallskip
    \textbf{Εντολές Maxima:}
    \texttt{load("eigen")}, \texttt{eigenvalues(A)},
    \texttt{eigenvectors(A)}, \texttt{charpoly(A,$\lambda$)},
    \texttt{nullspace(A-$\lambda$*I)}
  }%
  {https://nikmatz.github.io/ELADIC_GR/code/ch05/}%
  {qr_94ee53b0ae3e.png}

\subsection*{Δραστηριότητα Python}

\begin{pybox}{Python: Ιδιοτιμές — NumPy, Διαγωνοποίηση, Markov}%
             {https://nikmatz.github.io/ELADIC_GR/code/ch05/}%
             {qr_94ee53b0ae3e.png}
\label{py:ch5_eigenvalues}
\begin{enumerate}[label=(\alph*),itemsep=2pt]
  \item Με \texttt{sympy.Matrix.eigenvects()} βρείτε ακριβείς ιδιοτιμές
    και ιδιοδιανύσματα του $A$. Επαληθεύστε trace και det.
  \item Με \texttt{np.linalg.eig} επαληθεύστε αριθμητικά:
    $A\mathbf{v}_i = \lambda_i\mathbf{v}_i$.
  \item Κατασκευάστε $P$, $D$ και επαληθεύστε $PDP^{-1}=A$.
    Υπολογίστε $A^{10}$ μέσω $PD^{10}P^{-1}$ και συγκρίνετε με
    \texttt{np.linalg.matrix\_power}.
  \item Για τον πίνακα Markov
    $M=\bigl[\begin{smallmatrix}0.8&0.3\\0.2&0.7\end{smallmatrix}\bigr]$:
    βρείτε τη στάσιμη κατανομή $\pi$ (ιδιοδιάνυσμα $\lambda=1$)
    και επαληθεύστε $M\pi=\pi$.
  \item \textit{Προαιρετικό:} Οπτικοποιήστε σύγκλιση power iteration
    και Markov chain (Matplotlib).
\end{enumerate}
\medskip
\textbf{Βιβλιοθήκες / Εντολές:}
\texttt{np.linalg.eig()}, \texttt{sympy.Matrix.eigenvects()},
\texttt{sympy.Matrix.charpoly()},
\texttt{np.linalg.matrix\_power()}, \texttt{np.diag()}
\end{pybox}
